\chapter{Conclusão}
% ---

Ao final deste trabalho, conclui-se que foi possível reestruturar e
evoluir o ecossistema do Pró-Mamá, abrangendo a \emph{API}, o Painel
Administrativo e a infraestrutura de implantação. A reformulação não se
limitou à atualização tecnológica e buscou corrigir limitações do
sistema anterior, estabelecendo base mais adequada para manutenção,
continuidade e expansão da solução.

Entre os resultados alcançados, destaca-se a adoção de HTTPS com suporte
da Cloudflare, medida que corrigiu uma fragilidade importante de
segurança presente na versão anterior. Também se evidenciou a
reformulação do módulo de notificações, que deixou de depender de regras
embutidas no código e passou a operar por registros administráveis no
painel, tornando a configuração mais flexível pela equipe.

Outro resultado relevante foi a consolidação da infraestrutura como
código, centralizada em repositório específico para orquestração,
publicação e atualização do ambiente. Essa decisão fortaleceu a
reprodutibilidade da solução e tornou a implantação mais previsível,
criando condições para reaproveitar e escalar a arquitetura a outros
municípios interessados.

Do ponto de vista técnico, o trabalho contribuiu para a organização
interna do projeto, com padronização das rotas da \emph{API},
processamento assíncrono de notificações, documentação automatizada e
reorganização do painel administrativo. Em conjunto, essas decisões
favoreceram manutenção mais clara, separação de responsabilidades e
previsibilidade na evolução do sistema. Os testes automatizados do
\emph{back-end} reforçam esse cenário ao oferecer base verificável para
parte relevante das regras de negócio implementadas.

Apesar dos avanços obtidos, permanecem limitações relevantes. A ausência
de integração com bases oficiais do poder público mantém elevado o
esforço manual da equipe do NASF no cadastro de informações públicas e
reduz o foco nas interações com as mães atendidas pelo programa.

Outra limitação está na baixa divulgação e no baixo engajamento
institucional em torno da plataforma, tanto na prefeitura quanto no
IFRS, o que reduz o alcance prático de uma solução que já se encontra
funcional em produção, com aplicativo, painel e infraestrutura operando
de forma integrada.

Assim, conclui-se que o trabalho atingiu seu propósito ao entregar base
mais segura, administrável e sustentável para a continuidade do projeto.
Mais do que atualização tecnológica, a proposta preservou o valor social
do sistema e ampliou sua capacidade de apoiar, de forma digital, a
comunicação entre os serviços de saúde e as mães atendidas pelo
programa.

\section{Trabalhos Futuros}

Permanece um desafio de continuidade: a dimensão do projeto dificulta a
entrada de novos colaboradores e torna mais lento o entendimento do
domínio e da base de código. Um eixo promissor de evolução está no uso
de abordagens baseadas em inteligência artificial para apoiar o
\emph{onboarding} técnico e de negócio, como documentação orientada por
especificação (\emph{Specification-Driven Development}) e agentes
customizados em ferramentas de apoio ao desenvolvimento, a exemplo do
\emph{Copilot Chat}.

Outra frente relevante está na integração com bases oficiais do poder
público, de modo a reduzir cadastros manuais e concentrar a equipe do
NASF na interação com as mães. Esse avanço pode ampliar a eficiência do
uso da plataforma e deslocar o esforço institucional para atividades de
acompanhamento e orientação.

Também se mostra importante fortalecer ações de divulgação institucional
na prefeitura e no IFRS para ampliar uso, engajamento e colaboração
estudantil. Por fim, permanece como oportunidade ampliar os cenários de
validação, com mais testes integrados e outras formas de verificação
contínua do ecossistema.
